\documentclass[11pt]{article}
\usepackage[margin=1in]{geometry}
\usepackage{hyperref}
\usepackage{listings}
\usepackage{xcolor}
\usepackage{enumitem}
\usepackage{amsmath}
\usepackage{booktabs}

\emergencystretch=3em
\tolerance=2000
\sloppy

\lstset{
  basicstyle=\ttfamily\small,
  breaklines=true,
  columns=fullflexible,
  frame=single,
  backgroundcolor=\color{gray!10},
  showstringspaces=false
}

\title{\textbf{grep-dap}: Generative-AI assistance for validating and
       enriching the semantic metadata of remote Earth-science
       archives\\
       \large project summary and case study}
\author{P. Cornillon\\
\small Graduate School of Oceanography, University of Rhode Island}
\date{2026-05-18}

\begin{document}
\maketitle

\begin{abstract}
NASA's Earth-science archives expose petabytes of remote-sensing
and model data through OPeNDAP servers, but the \emph{semantic}
metadata that tells a user what those data \emph{mean}---units,
coordinates, provenance, processing history---varies enormously,
from fully CF-compliant down to none at all. The
\textbf{grep-dap} project explores how a generative-AI agent,
disciplined by the OPeNDAP protocol and by direct statistical
probing of the underlying data, can (a)~give archive curators a
prioritised, evidence-backed punch list of what is missing or
malformed in the metadata they are responsible for, and
(b)~produce ``usage'' artefacts that let downstream users and
their own AI agents work with the data even when the underlying
metadata is incomplete. This document summarises a six-step
proof-of-concept test case carried out against the URI OPeNDAP
server at \texttt{sst-aqua.gso.uri.edu} (the AVHRR/MODIS SST-Fronts
holdings), describes what the case shows is feasible, and frames
the resulting capability as the basis for a NASA-scale curator-
and-user support service.
\end{abstract}

\section{Motivation: a structural feature, not a defect, of OPeNDAP}

OPeNDAP (the DAP protocol) was designed with deliberately
asymmetric metadata: a rigid \emph{syntactic} structure (shapes,
dtypes, dimensions, groups) and a virtually free \emph{semantic}
structure (units, long names, standard names, provenance,
references) \cite{dap4}. The asymmetry lowered the barrier for
data providers to publish their holdings, at the cost of allowing
the semantic side to vary from complete to absent across servers.
In practice, this is exactly the distribution that one sees: a
single OPeNDAP server can host both CF-1.x--compliant ``flagship''
datasets and entirely silent derivative products in adjacent
directories. The disparity is visible \emph{within} the
17-directory URI server studied here.

A consequence of that variability is that an analyst who wants to
use a remote dataset must currently do, by hand, one of two
things:
\begin{itemize}[topsep=2pt]
  \item Trust the producer's intent and proceed, risking silent
    misinterpretation; or
  \item Reverse-engineer the dataset's semantics from the
    structural metadata, the file and directory naming, the
    statistical structure of the data itself, and any sibling
    datasets that happen to be better annotated.
\end{itemize}
The second path is what an experienced scientist does when they
encounter an unfamiliar dataset. It is also exactly the kind of
multi-modal, judgement-laden inspection at which contemporary
GenAI agents have become surprisingly capable.

\section{The grep-dap approach}\label{sec:approach}

\textbf{grep-dap} treats an OPeNDAP endpoint as a probeable
environment and an AI agent as a careful (sceptical, evidence-
demanding) external reader. The agent's reasoning is anchored
in three classes of evidence:

\begin{enumerate}[leftmargin=2em,topsep=2pt]
  \item \textbf{Catalog-level structure.} Directory names, file
    naming conventions, and version tags. A name like
    \texttt{AQUA\_MODIS\_orbit\_NNNNNN\_YYYYMMDDThhmmss\_L2\_SST-URI\_24-2.nc4}
    encodes platform, instrument, orbit number, start time,
    processing level, producer, and version without opening the
    file.
  \item \textbf{Structural metadata.} The DAP4 DMR (or DAP2
    DDS+DAS), parsed depth-aware so that every group, dimension,
    variable, dtype, shape, fill value, scale factor, and
    declared attribute is captured.
  \item \textbf{Data probes.} Small, deliberately chosen DAP
    hyperslab requests that test specific hypotheses (for example,
    ``is this variable in $^\circ$C?'', ``does lat run S--N or
    N--S?'', ``what is the effective spatial support of this
    derivative operator?''). These are the equivalent of a
    senior scientist asking ``does the field actually look like
    SST?'' before trusting anyone's metadata.
\end{enumerate}

The output of a grep-dap pass is not a single declarative answer
but a structured judgement that distinguishes
\emph{stated}, \emph{inferred}, and \emph{not safely guessable}
items. Inferences are themselves graded by confidence (statistical
support from data probes, naming-grammar regularity, sibling-
dataset consistency) so that downstream consumers know which
attributes are safe to act on automatically and which require
producer confirmation.

\section{The URI case study}\label{sec:case}

The 2026 test case ran a six-prompt sequence against the URI
OPeNDAP server at \texttt{https://sst-aqua.gso.uri.edu/opendap/}.
The server hosts 17 top-level containers; only two are publicly
readable (the remainder return Hyrax 403). All findings below
are evidence-tracked and reproducible from the scripts in the
project repository.

\subsection{Mapping the server (Prompt~1)}

The agent walked the root catalog, identified the server software
(Hyrax behind \texttt{nginx}, DAP2 + DAP4 both advertised),
classified all 17 directory names into thematic groups (sensor
sources \texttt{JAXA\_Orbits}, \texttt{RSS\_Orbits}, \texttt{MUR};
in-situ \texttt{iQuamBuoy}; matchups \texttt{matchups\_v2}
and \texttt{matchups\_v3};
time series \texttt{timeSeries\_v2}; gradient products
\texttt{gradients\_by\_period} and
\texttt{gradients\_by\_period\_5day}; orbit files
\texttt{SST\_Orbits}), and successfully read the two public
branches end-to-end. A key observation already at this step: a
forbidden directory's \emph{name} is still useful semantic signal
that catalog-level inference should not discard.

\subsection{COARDS audit of a well-annotated dataset (Prompt~2)}

For the CF-1.5--labelled \texttt{SST\_Orbits/} branch, the agent
parsed the DMR of a representative file into a complete inventory
(3 groups, 3 + 6 + 3 dimensions, 44 atomic variables, 22
root-group attributes) and scored each variable against a
deliberately strict reading of COARDS. Of 44 variables, 21 were
COARDS-complete and 23 were not, with three recurring issues that
accounted for almost all of the gap:

\begin{itemize}[topsep=2pt]
  \item Non-UDUNITS unit strings (\texttt{units="C"},
    \texttt{units="C/km"}) on 15 variables; one substitution per
    variable resolves them.
  \item Three time variables (\texttt{DateTime}, two
    \texttt{*\_time} in the \texttt{/contributing\_granules/}
    group) declaring \texttt{units="seconds"} without an epoch
    reference.
  \item Dimensionless quantities (region indices, granule indices)
    with no \texttt{units} attribute at all.
\end{itemize}

The agent additionally identified two factually wrong unit
declarations buried in a sub-group (\texttt{L2eqa\_AMSR\_E\_wind\_speed}
declared \texttt{units="mm"}; \texttt{L2eqa\_MODIS\_num\_SST}
declared \texttt{units="C"} on a pixel count), each by comparing
the declared units with the packed value range and with the
sibling variables' units. These would have been hard for a
human-only review to find at scale.

For each affected variable the audit emitted a proposed value
and the evidence that supports it; for each variable, the audit
\emph{also} suggested non-COARDS attributes (\texttt{coordinates},
\texttt{ancillary\_variables}, \texttt{cell\_methods},
\texttt{flag\_values}/\texttt{flag\_meanings}) that would
significantly aid interpretation. Two factual errors in an
earlier exploratory pass (Prompt~1) were caught and explicitly
corrected as part of the audit.

\subsection{Recovering semantics for a zero-metadata dataset (Prompt~3)}

The \texttt{gradients\_by\_period/} branch (270 monthly NetCDF
files of statistics over MODIS-Aqua SST and SST gradients,
2003--2024, on a $1^\circ \times 1^\circ$ global grid) carries
\emph{zero} attributes: no \texttt{units}, no
\texttt{long\_name}, no \texttt{\_FillValue}, no coordinate
variables, no global attributes. This is the OPeNDAP free-for-all
that motivates grep-dap.

The agent recovered the dataset's semantics by combining three
independent lines of evidence:

\begin{enumerate}[topsep=2pt]
  \item \textbf{Variable-name grammar} (e.g.\ the regular structure
    \texttt{\{day,night\}\_sum\_\{SST,\dots\}\_squared} implying
    a sum + sum-of-squares + count bookkeeping scheme).
  \item \textbf{Sibling-dataset inheritance} (the sibling
    \texttt{SST\_Orbits/} declares the gradient variables in
    $^\circ$C/km; the gradient sums in this dataset inherit
    those units).
  \item \textbf{Direct geographic probes} (classifying six known-
    land and six known-ocean points against the
    \texttt{day\_pixel\_count} field. Only one of the four
    candidate lat/lon conventions gets all 12 correct, giving
    \texttt{lat[k] = -89.5 + k} (south-to-north) and
    \texttt{lon[k] = -179.5 + k} ($-180\to+180$) unambiguously).
\end{enumerate}

The audit was organised as a $3 \times 2$ matrix:
inference confidence (\emph{very confident} / \emph{guessed} /
\emph{will not guess}) crossed with importance (\emph{essential} /
\emph{nice-to-have}). This explicitly classifies which inferred
attributes the agent is willing to write back into a re-published
file and which it refuses to fabricate.

\subsection{Estimating the physical scale of a derived product (Prompt~4)}

For the gradient product, the agent estimated the \emph{spatial
range over which a single gradient value was computed} (the
support of the differentiation operator in the upstream L2 swath)
by three independent methods:
\begin{itemize}[topsep=2pt]
  \item Pixel-count ratios (the \texttt{at\_pixel\_count} /
    \texttt{pixel\_count} ratio is sensitive to scan-boundary
    masking of the gradient stencil, fixing the operator half-
    width at $\sim$2 pixels).
  \item Magnitude-RMSE matching against centered differences of
    the masked SST field at multiple scales (RMSE minimum at $h=3$
    pixels, support $\sim$6\,km).
  \item Autocorrelation of the stored gradient field
    (decorrelation length $\sim$2--3\,km).
\end{itemize}
All three converge on $5 \pm 2$\,km, an estimate that no part of
the file's declared metadata supplies but that is essential for
anyone using the data in a quantitative model. The agent recorded
the estimate along with what the analysis explicitly cannot
resolve (the precise functional form of the operator; whether
upstream SST smoothing was applied in addition to masking).

\subsection{Producing user-facing artefacts (Prompts~5--6)}

Three classes of artefact were produced from the audits:

\begin{itemize}[topsep=2pt]
  \item \textbf{User-facing dataset descriptions.} Short, accessible
    overviews of each dataset for a working scientist deciding
    whether the data fit a need, including a 15-line code snippet
    for a subset and an explicit ``honest limitations'' list.
  \item \textbf{Agent-facing usage files} (\texttt{usage\_sst.md},
    \texttt{usage\_gradient\_SST.md}). Self-contained Markdown
    guides intended to be handed to a downstream AI agent or
    pipeline. They are written in a form an LLM agent can act on
    directly: server endpoint, access pattern, variable inventory
    with all inferences marked, code recipes for the common cases,
    and a numbered list of pitfalls. The headline code examples
    were verified against the live server in a sanity-test run.
  \item \textbf{Curator-facing letters} (one per dataset). Short
    (1--2 pages), prioritised punch-lists of metadata gaps,
    written in the language a producing scientist would expect:
    explicit attribute name, proposed value, evidence, and a
    link back to the deeper audit. The CF-compliant dataset's
    letter is mostly small unit-string substitutions; the
    zero-metadata dataset's letter splits into ``what can be
    confidently encoded'' and ``what only the producer can
    supply'' (date\_created, processing history, exact upstream
    algorithm).
\end{itemize}

\section{Capabilities the test case demonstrates}\label{sec:caps}

The URI test case exercises six capabilities that an automated
service would need at scale:

\begin{enumerate}[leftmargin=2em,topsep=2pt]
  \item \textbf{Catalog walking with permission-aware harvesting.}
    A grep-dap agent must extract signal from directory names it
    cannot read, and must not penalise a curator for having
    private holdings on the same server.
  \item \textbf{Depth-aware DMR parsing.} Hyrax/THREDDS DMR XML
    has nested groups; a naive flat-regex parser (as I used in
    Prompt~1) silently misses both global attributes and the
    full per-variable attribute set. A first-class agent uses a
    proper namespace- and depth-aware parser, and re-runs prior
    inferences when corrected.
  \item \textbf{Inference under confidence.} Each proposed
    attribute is tagged with the type of evidence that supports
    it; the agent refuses to fabricate items it cannot defend
    (\texttt{date\_created}, \texttt{history}, the exact form of
    a black-box operator).
  \item \textbf{Cross-dataset reasoning.} The agent uses better-
    annotated sibling datasets on the same server (here:
    \texttt{SST\_Orbits/} informing \texttt{gradients\_by\_period/})
    as a primary inference source. Many archive layouts share
    this property.
  \item \textbf{Quantitative inference from data probes.} The
    spatial-scale estimation in Prompt~4 is a small example of a
    much larger class: physically meaningful quantities
    (resolution, decorrelation lengths, noise floor, latency,
    diurnal/seasonal cycle phase) that can be estimated from
    cheap DAP hyperslab requests.
  \item \textbf{Producing distinct artefacts for distinct
    audiences.} A curator wants a punch list; a user wants a
    pre-flight summary plus working code; an LLM agent wants a
    machine-readable usage file. The same audit drives all three.
\end{enumerate}

\section{Implications for a NASA-scale service}\label{sec:nasa}

NASA's Earthdata catalog spans dozens of DAACs and thousands of
collections served via OPeNDAP-compatible endpoints (Hyrax,
THREDDS, ERDDAP). The variability documented within one URI
server is, anecdotally, mirrored across the system. A grep-dap--
like service offered to curators and users at NASA scale could:

\begin{itemize}[topsep=2pt]
  \item \textbf{Continuously audit} every collection's DMR against
    CF, COARDS, and ACDD requirements, with prioritisation tuned
    to the discipline (e.g.\ flag missing time epochs and
    coordinate variables harder than missing \texttt{keywords}).
  \item \textbf{Cross-link related datasets.} Inheritance from
    well-annotated to thinly-annotated companions is a powerful
    inference channel that no single-file validator currently
    uses. NASA holdings, with their explicit collection-of-
    granules structure, are particularly amenable to this.
  \item \textbf{Estimate physical quantities by data probing.}
    For derived products, the producer's chosen resolution,
    operator support, and noise floor are often not
    documented---and yet are essential for the user to decide
    whether the data fit a science question. Cheap, principled
    probing can recover these.
  \item \textbf{Generate dual-audience artefacts.} A weekly
    pull-request--style notification to curators (``these three
    attributes were either malformed or missing on collection X
    last week; here are the proposed fixes; click `Accept' to
    push the change to the producing script'') paired with
    public, machine-readable usage documents (regenerated
    automatically as collections change) would close the loop.
  \item \textbf{Provide an evidence trail.} Every proposed change
    or inference is tied to the structural feature, naming
    grammar item, or data probe that supports it; nothing is
    fabricated. That property is the project's main scientific
    discipline and a precondition for trust in an automated
    metadata system.
\end{itemize}

\section{Limitations and open questions}\label{sec:limits}

The URI case study is deliberately small in scope, and the
generalisations above must be tested rather than assumed.

\begin{itemize}[topsep=2pt]
  \item \textbf{Sample size.} A single 2-readable-directory server
    is not a representative sample of NASA archives. Scaling
    requires running grep-dap across a stratified sample of DAACs
    (PO.DAAC, GES~DISC, NSIDC, ASDC, LP~DAAC, etc.) and
    quantifying the audit error rate against curator review.
  \item \textbf{Agent reliability.} Two factual errors in Prompt~1
    of the case study (a missed root-attribute set and a missed
    scale factor) were caught and corrected in Prompt~2. A
    deployable service must have automatic correction loops and
    independent re-runs as a built-in feature.
  \item \textbf{Producer overhead.} The curator-facing letter is
    intentionally short, but pushing it into producers' actual
    workflows (notebooks, processing scripts, configuration
    files) requires integration with their tooling. Designing
    that interface --- and especially the ``one-click accept''
    pull-request workflow --- is itself non-trivial.
  \item \textbf{Authentication.} 14 of 17 directories on the URI
    server are 403 from a public client. NASA Earthdata's
    \texttt{urs.earthdata.nasa.gov} login provides an obvious
    authentication path for a service that operates on behalf of
    the archive owner; an external-only auditor would not see
    those collections at all.
  \item \textbf{Failure modes of GenAI on metadata.} An agent that
    cheerfully fabricates a plausible \texttt{date\_created} or
    \texttt{processing\_history} is worse than no audit at all.
    The strict confidence-classification of inferences is one
    line of defence; explicit refusal-to-guess for unsupported
    fields is another. Both need empirical evaluation.
\end{itemize}

\section{Proposed next steps}\label{sec:next}

Three immediate extensions of the URI proof-of-concept would
position the work for a NASA-scale follow-on:

\begin{enumerate}[topsep=2pt,leftmargin=2em]
  \item \textbf{Multi-DAAC pilot.} Run grep-dap against one
    representative collection from each of three contrasting
    DAACs (e.g.\ PO.DAAC L2 SST, GES~DISC GPM L3, NSIDC sea-ice
    L4). Compare the audit output against curator review and
    quantify precision/recall on (i) flagged COARDS issues and
    (ii) proposed attribute values.
  \item \textbf{Producer-loop integration.} Build a prototype
    that issues GitHub-style pull requests against the producing
    repositories for a small number of cooperating curators,
    closing the loop from external audit to in-source attribute
    change.
  \item \textbf{User-facing usage-doc service.} Auto-generate
    \texttt{usage\_*.md} files for every audited collection and
    expose them through a simple discovery interface. Measure
    user-task success rates (e.g.\ time to first working
    subsetting query) with and without the usage doc.
\end{enumerate}

\section{Reproducibility}\label{sec:repro}

All artefacts of the URI case study are checked into the
\texttt{grep-dap} repository:

\begin{itemize}[topsep=2pt]
  \item \texttt{docs/uri\_test\_case\_\{1\dots5\}.tex} --- the
    audit reports for each prompt.
  \item \texttt{docs/test\_case\_uri\_log\_\{1\dots6\}.tex} ---
    the timestamped working logs.
  \item \texttt{docs/curator\_report\_sst\_orbits.tex},
    \texttt{docs/curator\_report\_gradients\_by\_period.tex} ---
    the two curator letters.
  \item \texttt{docs/usage\_sst.md},
    \texttt{docs/usage\_gradient\_SST.md} --- the agent-facing
    usage documents.
  \item \texttt{scripts/*.py} --- the probe scripts (catalog
    walk, DMR parser, COARDS scorer, geographic-orientation
    test, gradient-scale estimators, usage sanity tests).
\end{itemize}

Each artefact links to the scripts and inputs that produced it,
so any claim in the audits can be re-verified by re-running
against the live server.

\begin{thebibliography}{9}
\bibitem{dap4} Unidata and OPeNDAP, Inc., ``DAP4 Specification,''
  2016. \url{https://opendap.github.io/dap4-specification/DAP4.html}
\bibitem{cf18} CF Conventions and Metadata, version 1.8,
  \url{https://cfconventions.org/}
\bibitem{acdd} Unidata, ``Attribute Convention for Data
  Discovery,'' v1.3,
  \url{https://wiki.esipfed.org/Attribute_Convention_for_Data_Discovery_1-3}
\bibitem{earthdata} NASA Earthdata,
  \url{https://earthdata.nasa.gov/}
\end{thebibliography}

\end{document}
